\documentclass[a4paper,12pt]{article}
\usepackage[T2A]{fontenc}
\usepackage[utf8]{inputenc}
\usepackage[russian]{babel}
\usepackage{amsmath,amssymb,mathtools}
\usepackage{icomma}
\usepackage{geometry}
\geometry{left=23mm,right=23mm,top=24mm,bottom=24mm}
\usepackage{microtype}
\usepackage{tikz}
\usetikzlibrary{arrows.meta,positioning,shapes.geometric,calc}
\usepackage{tabularx,booktabs,longtable}
\usepackage{enumitem}
\usepackage{fancyhdr}
\usepackage{setspace}
\usepackage{xcolor}
\usepackage{titlesec}
\usepackage[hidelinks]{hyperref}
\definecolor{accent}{HTML}{008080}
\definecolor{darkaccent}{HTML}{00575B}
\definecolor{textgray}{HTML}{333333}
\definecolor{softgray}{HTML}{667078}
\color{textgray}
\setlength{\parindent}{0pt}
\setlength{\parskip}{5pt}
\setlist[enumerate]{itemsep=7pt,topsep=6pt,leftmargin=7mm}
\setlist[itemize]{itemsep=4pt,topsep=5pt,leftmargin=7mm}
\setstretch{1.19}
\titleformat{\section}{\Large\color{darkaccent}}{\thesection.}{0.55em}{}[\vspace{3pt}\titlerule]
\titlespacing*{\section}{0pt}{15pt}{10pt}
\titleformat{\subsection}{\normalsize\color{darkaccent}\bfseries}{\thesubsection.}{0.5em}{}
\titlespacing*{\subsection}{0pt}{11pt}{4pt}
\pagestyle{fancy}
\fancyhf{}
\fancyhead[L]{\footnotesize Повторение алгебры и метод интервалов}
\fancyhead[R]{\footnotesize Володя Хачатурян}
\fancyfoot[C]{\footnotesize\thepage}
\renewcommand{\headrulewidth}{0.35pt}
\renewcommand{\headrule}{\hbox to\headwidth{\color{accent}\leaders\hrule height \headrulewidth\hfill}}
\setlength{\headheight}{23pt}
\newcommand{\idea}[1]{\textcolor{darkaccent}{\textbf{Главная идея.}} #1}
\newcommand{\attention}[1]{\textcolor{darkaccent}{\textbf{Проверьте себя.}} #1}
\newcommand{\ex}[1]{\textcolor{darkaccent}{\textbf{Пример.}} #1}
\begin{document}
\thispagestyle{empty}
\vspace*{31mm}
\begin{center}
{\large\color{darkaccent} МАТЕМАТИКА}\par
\vspace{17mm}
{\LARGE Чек-лист}\par
\vspace{11mm}
{\LARGE\color{accent} Алгебраические преобразования,\\[2mm] системы и метод интервалов}\par
\vspace{13mm}
{\normalsize Повторение базовых приёмов и подготовка к изучению неравенств}\par
\vspace{32mm}
{\large Володя Хачатурян}\par
\vspace{13mm}
{\normalsize 23 ноября 2026 г.}\par
\vfill
{\small Лёвин Артём Александрович\\ эксклюзивно для Володи Хачатуряна}\par
\vspace{18mm}
\begin{tikzpicture}[x=1cm,y=1cm]
\draw[accent,line width=.85pt] (0,0) .. controls (1.0,.47) and (2.0,-.48) .. (3.0,0)
.. controls (4.0,.47) and (5.0,-.48) .. (6.0,0);
\end{tikzpicture}
\end{center}
\clearpage

\section{Линейные уравнения и пропорции}
\idea{После каждого преобразования проверяйте, сохраняется ли равенство.}

Если \(\frac{a}{b}=c\) и \(b\ne0\), \(c\ne0\), то \(b=\frac{a}{c}\).
\[
\frac{2}{x}=10\quad\Longrightarrow\quad 2=10x
\quad\Longrightarrow\quad x=\frac15=0{,}2.
\]
\ex{При переносе слагаемого в другую часть уравнения его знак меняется:}
\[
2-x=10\quad\Longrightarrow\quad -x=8\quad\Longrightarrow\quad x=-8.
\]

При равенстве двух дробей с ненулевыми знаменателями удобно применить перекрёстное умножение. \textbf{Произведения с суммой или разностью записывайте в скобках.}
\[
\frac{x}{3}=\frac{2x+5}{4}
\Longrightarrow 4x=3(2x+5)
\Longrightarrow 4x=6x+15
\Longrightarrow x=-\frac{15}{2}.
\]
\attention{Проверка подстановкой позволяет обнаружить ошибочный перенос или потерянную скобку.}

\section{Линейные неравенства и промежутки}
При умножении или делении обеих частей неравенства на \emph{отрицательное число} знак \(>,<,\ge,\le\) меняется на противоположный.
\[
4x+3\ge7x+5\quad\Longrightarrow\quad -3x\ge2
\quad\Longrightarrow\quad \boxed{x\le-\frac23}.
\]
Для \(x\le-\frac23\) граничная точка входит в решение. На числовой оси она закрашена, а в записи промежутка справа стоит квадратная скобка:
\[
x\in\left(-\infty;-\frac23\right].
\]
\begin{center}
\begin{tikzpicture}[scale=.95]
\draw[-{Stealth[length=3mm]},gray] (-4.8,0)--(4.5,0);
\draw[accent,line width=2.2pt] (-4.65,0)--(.45,0);
\fill[accent] (.45,0) circle (2.5pt);
\node[below=5pt] at (.45,0) {$-\frac23$};
\end{tikzpicture}
\end{center}
\attention{Для строгого неравенства \(<\) или \(>\) граничная точка не включается; бесконечности всегда записывают с круглыми скобками.}

\section{Квадратные уравнения}
Для \(ax^2+bx+c=0\), где \(a\ne0\), дискриминант и корни:
\[
D=b^2-4ac,\qquad x_{1,2}=\frac{-b\pm\sqrt D}{2a}.
\]
\begin{center}
\begin{tabular}{@{}ll@{}}
\toprule
\(D>0\) & два различных действительных корня;\\
\(D=0\) & один действительный корень;\\
\(D<0\) & действительных корней нет.\\
\bottomrule
\end{tabular}
\end{center}
\ex{Для \(x^2-3x-4=0\):}
\[
D=(-3)^2-4\cdot1\cdot(-4)=25,
\quad x_1=\frac{3+5}{2}=4,
\quad x_2=\frac{3-5}{2}=-1.
\]
Если свободного члена нет, часто удобнее вынести \(x\):
\[
x^2-3x=x(x-3)=0\quad\Longrightarrow\quad x=0\ \text{или}\ x=3.
\]
Для \(x^2+4=0\) действительных корней нет, поскольку \(x^2\ge0\) и \(x^2+4>0\).

\section{Вынесение общего множителя}
\idea{Найдите общий числовой и буквенный множитель каждого слагаемого, затем вынесите его за скобки.}
\[
a^2b+ab^2=ab(a+b),\qquad
x^2-3x=x(x-3).
\]
\attention{После разложения полезно раскрыть скобки обратно: должно получиться исходное выражение.}

\section{Формулы сокращённого умножения}
Три формулы из занятия различаются по \emph{расположению скобок и степени}:
\begin{align*}
(a-b)^2&=a^2-2ab+b^2&&\text{(квадрат разности)},\\
(a+b)^2&=a^2+2ab+b^2&&\text{(квадрат суммы)},\\
a^2-b^2&=(a-b)(a+b)&&\text{(разность квадратов)}.
\end{align*}
\ex{Вместо \(a\) и \(b\) подставляем \emph{целые выражения}:}
\begin{align*}
(3x-5y)^2
&=(3x)^2-2\cdot(3x)\cdot(5y)+(5y)^2\\
&=9x^2-30xy+25y^2,\\[2pt]
(4t+7k)^2
&=(4t)^2+2\cdot(4t)\cdot(7k)+(7k)^2\\
&=16t^2+56tk+49k^2,\\[2pt]
4t^2-9k^2&=(2t-3k)(2t+3k).
\end{align*}
\idea{Любое положительное число \(c\) представимо как \((\sqrt c)^2\). Поэтому}
\[
x^2-3=(x-\sqrt3)(x+\sqrt3),
\quad x^2-3=0\Longrightarrow x=\pm\sqrt3.
\]
\attention{В общем случае \((a-b)^2\ne a^2-b^2\). В квадрате разности обязательно есть удвоенное произведение \(-2ab\).}

\section{Алгебраические дроби и «улыбочка»}
Для сложения или вычитания дробей сначала приводим их к общему знаменателю:
\[
\frac{A}{B}-\frac{C}{D}=\frac{AD-CB}{BD},\qquad B\ne0,\ D\ne0.
\]
\ex{Пример из занятия. ОДЗ: \(x\ne-3\), \(x\ne3\).}
\begin{align*}
\frac{x-3}{x+3}-\frac{x+3}{x-3}
&=\frac{(x-3)^2-(x+3)^2}{(x+3)(x-3)}\\
&=\frac{(x^2-6x+9)-(x^2+6x+9)}{x^2-9}\\
&=\boxed{\frac{-12x}{x^2-9}}.
\end{align*}
\attention{Скобки обязательны вокруг всего вычитаемого числителя. Сокращать можно только общие \emph{множители} числителя и знаменателя, сохраняя ОДЗ.}

\section{Системы неравенств: пересечение}
\idea{Каждое неравенство задаёт множество чисел. Решение системы --- их общая часть.}
\[
\begin{cases}
x\ge-5,\\
x\le10
\end{cases}
\quad\Longrightarrow\quad x\in[-5;10].
\]
\ex{Два линейных неравенства:}
\[
\begin{cases}
4x+5\ge10x-7,\\
x-1\ge7x-5
\end{cases}
\Longleftrightarrow
\begin{cases}
x\le2,\\
x\le\frac23
\end{cases}
\Longrightarrow\boxed{x\in\left(-\infty;\frac23\right]}.
\]
Слово \emph{система} означает «одновременно»: берём пересечение, даже если одно условие сильнее другого.

\section{Системы уравнений: подстановка}
\idea{Выразите одну переменную через другую, подставьте во второе уравнение, затем найдите оставшуюся переменную.}
\[
\begin{cases}x+y=10,\\2x-3y=15\end{cases}
\quad\Longrightarrow\quad x=10-y.
\]
\[
2(10-y)-3y=15\Longrightarrow20-5y=15
\Longrightarrow y=1\Longrightarrow x=9.
\]
\[
\boxed{(x;y)=(9;1)}.
\]
\attention{Проверяйте \emph{оба} исходных уравнения: \(9+1=10\), \(2\cdot9-3\cdot1=15\).}

\clearpage
\section{Начало метода интервалов}
Неравенство \(f(x)\ge0\) требует найти те значения \(x\), при которых выражение \(f(x)\) \emph{неотрицательно}. Для \(f(x)\le0\) ищем неположительные значения. Нуль включается в обоих случаях.

\ex{Решим \(x^2-x-6\ge0\).}
\[
x^2-x-6=(x+2)(x-3).
\]
Нули: \(-2\) и \(3\). Они разбивают прямую на три промежутка. Проверяем знак в каждом:
\begin{center}
\begin{tabular}{@{}cccc@{}}
\toprule
Промежуток & \(x\) для проверки & \(f(x)\) & знак\\
\midrule
\(( -\infty;-2)\) & \(-3\) & \(9+3-6=6\) & \(+\)\\
\((-2;3)\) & \(0\) & \(-6\) & \(-\)\\
\((3;+\infty)\) & \(4\) & \(16-4-6=6\) & \(+\)\\
\bottomrule
\end{tabular}
\end{center}
\begin{center}
\begin{tikzpicture}[x=.94cm,y=1cm]
\draw[-{Stealth[length=3mm]},gray] (-5,0)--(5.3,0);
\draw[accent,line width=2.2pt] (-4.85,0)--(-1.8,0);
\draw[accent,line width=2.2pt] (1.8,0)--(5.03,0);
\fill[accent] (-1.8,0) circle (2.4pt);
\fill[accent] (1.8,0) circle (2.4pt);
\node[below=5pt] at (-1.8,0) {$-2$};
\node[below=5pt] at (1.8,0) {$3$};
\node[above=5pt] at (-3.35,0) {$+$};
\node[above=5pt] at (0,0) {$-$};
\node[above=5pt] at (3.4,0) {$+$};
\end{tikzpicture}
\end{center}
\[
\boxed{x\in(-\infty;-2]\cup[3;+\infty)}.
\]
\ex{Для \(x^2-3x-4\le0\) корни равны \(-1\) и \(4\). Внешние интервалы имеют знак «плюс», внутренний --- «минус», поэтому \(x\in[-1;4]\).}

\textcolor{darkaccent}{\textbf{Важное уточнение.}} Знаки обязательно меняются при прохождении через корень \emph{нечётной кратности}. Если корень имеет чётную кратность, знак может сохраняться. Для рассмотренных квадратных трёхчленов с двумя различными корнями чередование работает.

\clearpage
\thispagestyle{fancy}
\begin{center}
{\Large\color{darkaccent} Алгоритм: метод интервалов для многочлена}\par
\vspace{15mm}
\tikzset{
 proc/.style={draw=darkaccent,rounded corners=3pt,text width=9.2cm,minimum height=13mm,align=center,inner sep=4pt},
 decision/.style={draw=darkaccent,diamond,aspect=3,text width=6.3cm,align=center,inner sep=1pt},
 term/.style={draw=accent,ellipse,text width=6.2cm,minimum height=11mm,align=center},
 ar/.style={-{Stealth[length=2.8mm]},line width=.9pt,draw=darkaccent}
}
\begin{tikzpicture}[node distance=9mm]
\node[term] (start) {Начало: определить знак \(f(x)\)};
\node[proc,below=of start] (root) {Решить \(f(x)=0\); отметить все действительные корни на оси};
\node[proc,below=of root] (test) {В каждом промежутке выбрать удобную точку\\ и вычислить знак \(f(x)\)};
\node[proc,below=of test] (choose) {Выделить промежутки с нужным знаком:\\ «плюс» для \(>0,\ge0\); «минус» для \(<0,\le0\)};
\node[decision,below=13mm of choose] (strict) {Неравенство строгое?};
\node[proc,text width=4.3cm,below left=16mm and 12mm of strict] (yes) {Корни в ответ\\не включать};
\node[proc,text width=4.3cm,below right=16mm and 12mm of strict] (no) {Корни включить,\\если \(f(x)=0\)};
\node[term,below=37mm of strict] (finish) {Записать объединение промежутков};
\draw[ar] (start)--(root);
\draw[ar] (root)--(test);
\draw[ar] (test)--(choose);
\draw[ar] (choose)--(strict);
\draw[ar] (strict.south west) -| node[pos=.3,above left,font=\footnotesize]{да} (yes.north);
\draw[ar] (strict.south east) -| node[pos=.3,above right,font=\footnotesize]{нет} (no.north);
\draw[ar] (yes.south) |- (finish.west);
\draw[ar] (no.south) |- (finish.east);
\end{tikzpicture}
\end{center}
\clearpage

\section*{Тренировочный блок}
\textcolor{darkaccent}{\textbf{Путь А --- 10 заданий от простого к сложному.}} Решения запишите последовательно. Обратите внимание на ОДЗ, знаки неравенств и скобки.
\begin{enumerate}
\item Решите уравнение: \(\displaystyle\frac{x+2}{4}=\frac{3x-1}{5}\).
\item Решите неравенство. Ответ запишите промежутком: \(6x+1\ge2x+13\).
\item Решите квадратное уравнение: \(x^2-5x-14=0\).
\item Вынесите общий множитель: \(6ab^2+9a^2b\).
\item Раскройте скобки: \((2x-3y)^2\).
\item Разложите на множители: \(25p^2-16q^2\).
\item Упростите, указав ОДЗ:
\[\frac{x-2}{x+2}-\frac{x+2}{x-2}.\]
\item Решите систему неравенств:
\[\begin{cases}4x+1\ge2x-5,\\5x-4\le2x+8.\end{cases}\]
\item Решите методом подстановки систему уравнений:
\[\begin{cases}x+y=7,\\2x-y=5.\end{cases}\]
\item Решите методом интервалов: \(x^2-2x-8\le0\). Постройте схему знаков.
\end{enumerate}
\vspace{7mm}
\attention{Дополнительное задание на неделю: присылайте фотографии классных и домашних работ по школьной математике для согласования дальнейшей подготовки.}
\clearpage
\section*{Ответы}
\vspace{4mm}
\begin{center}
\renewcommand{\arraystretch}{1.5}
\begin{tabularx}{\linewidth}{@{}p{13mm}X@{}}
\toprule
№ & Ответ \\
\midrule
1 & \(x=2\).\\
2 & \(x\in[3;+\infty)\).\\
3 & \(x_1=-2,\quad x_2=7\).\\
4 & \(3ab(2b+3a)\).\\
5 & \(4x^2-12xy+9y^2\).\\
6 & \((5p-4q)(5p+4q)\).\\
7 & \(\displaystyle-\frac{8x}{x^2-4}\), при \(x\ne-2\), \(x\ne2\).\\
8 & \(x\in[-3;4]\).\\
9 & \((x;y)=(4;3)\).\\
10 & \(x\in[-2;4]\).\\
\bottomrule
\end{tabularx}
\end{center}
\vfill
\begin{center}
\small\color{softgray}Лёвин Артём Александрович эксклюзивно для Володи Хачатуряна
\end{center}
\end{document}
